% synthesis.tex -- revised discussion and conclusions for the updated report.
% Every number here is taken from the section it cites; nothing is introduced for the first time.

\section{Revised Discussion and Conclusions}
\label{sec:synthesis}

Sections~\ref{sec:orig-discussion} and~\ref{sec:orig-summary} are the discussion and summary as they stood at
Checkpoint~2. This section replaces them for the present state of the project. It first sets the main
conclusions of Checkpoint~2 against what the repository now shows, then draws the new studies together.

\subsection{The Checkpoint 2 conclusions, revisited}
\label{sec:synthesis-revisited}

Table~\ref{tab:synthesis-revisited} lists the conclusions of Checkpoint~2 and their present standing. The
complete list of corrected statements, including individual table entries, is
Table~\ref{tab:core-corrections}.

\begingroup
\small
\setlength{\tabcolsep}{4pt}
\begin{longtable}{@{}p{0.27\linewidth}p{0.50\linewidth}p{0.17\linewidth}@{}}
\caption{The main conclusions of the Checkpoint~2 report and their standing at commit \texttt{61d18d8}.}
\label{tab:synthesis-revisited}\\
\toprule
Checkpoint 2 conclusion & Standing now & Evidence \\
\midrule
\endfirsthead
\toprule
Checkpoint 2 conclusion & Standing now & Evidence \\
\midrule
\endhead
Zschokke's construction sequence can be written as an executable HTN
(Sections~\ref{sec:orig-prelim}--\ref{sec:orig-energy})
& Stands. The model now also generates explicit spiral geometry (logarithmic auxiliary spiral, Archimedean
capture spiral), and exists in a configurable form with 56 passing tests.
& Sections~\ref{sec:core-plans}, \ref{sec:core-spiral}, \ref{sec:core-tests} \\
\addlinespace
Early-stage variability is reproduced by backtracking over three proto-hub strategies
(Sections~\ref{sec:orig-htn}, \ref{sec:orig-protohub})
& Encoded, but not exercised. The planner as configured is greedy and does not backtrack, and only the
\textit{via-walk} strategy has ever appeared in a generated plan.
& Sections~\ref{sec:recovery-strategy}, \ref{sec:core-config} \\
\addlinespace
The default plan has 70 actions and costs 600.1\,cm (Section~\ref{sec:orig-vis})
& The plan with that cost has 71 actions and came from the configurable package as it then stood. The current
plan of that package has 105 actions and 701.6\,cm; the original package's plan has 128 actions and
751.6\,cm.
& Section~\ref{sec:core-plans} \\
\addlinespace
The HTN agent uses about 14.3 times less locomotory energy than a random walker
(Section~\ref{sec:orig-comparison})
& The order-of-magnitude gap stands: over 1{,}000 trials per condition the ratio of mean energies is
13.2--14.7 and of medians 10.3--11.5. The particular figures of Table~\ref{tab:orig-comparison} do not: its
HTN run had anchor removal switched on, and its baseline column cannot be reproduced.
& Sections~\ref{sec:core-anchor}, \ref{sec:core-compare} \\
\addlinespace
The HTN agent stays more reliable than the baseline as anchors are removed
(Sections~\ref{sec:orig-anchor-removal}, \ref{sec:orig-discussion})
& Not supported. The two success rates use different criteria; the HTN rate equals the probability that at
least one removable anchor survives, and the baseline's rate is higher from 60\,\% removal onwards.
& Section~\ref{sec:core-anchor} \\
\addlinespace
Wind is not yet swept (Table~\ref{tab:orig-roadmap})
& A wind sweep exists. In this model wind changes the cost that is reported and nothing else: success is
unchanged at every noise level for both agents.
& Section~\ref{sec:core-wind} \\
\addlinespace
The random-walker comparison isolates the effect of hierarchy (Section~\ref{sec:orig-baseline})
& It does not. Under a stricter flat baseline, the same action sequence without compound tasks, the HTN
agent and the flat agent perform identically in the moving-support benchmark (claim C06).
& Sections~\ref{sec:core-compare}, \ref{sec:mulder-claims} \\
\addlinespace
The task hierarchy reaches an efficient plan at a fraction of the search cost of a flat-state search
(Section~\ref{sec:orig-discussion})
& Not measured. The repository's flat-state planner (A*) solves a much weaker goal at the commit where it
works, and returns no plan at the current commit because of a defect in its heuristic.
& Sections~\ref{sec:astar_fap-effort}, \ref{sec:astar_fap-compare} \\
\bottomrule
\end{longtable}
\endgroup

\subsection{What the work since Checkpoint 2 shows}
\label{sec:synthesis-new}

\paragraph{The model itself.}
The HTN formalisation has become a configurable, tested and instrumented package whose plan contains real
spiral geometry rather than a placeholder polygon (Section~\ref{sec:core}). Auditing it against its own
committed results also showed where the earlier account was loose: which plan the published numbers
described, what ``success'' meant for each agent, and which stochastic settings were active.

\paragraph{Alternative formulations.}
Two non-HTN formulations exist beside the HTN models (Section~\ref{sec:astar_fap}). The fixed-action-pattern
model is a method library with exactly one method per task: it expands to a fixed 61-action sequence and has
no alternative when a precondition fails. The A* formulation searches a flat state space; as committed it
finds a 26-action plan for a goal that only counts threads, and since a later change to a shared distance
function its heuristic evaluates to zero and the search finds no plan within its budget. The four planners
are therefore described side by side but not ranked (claim C17).

\paragraph{Construction on moving supports.}
The moving-support benchmark (Section~\ref{sec:mulder}) gives the clearest result of the project so far, and
it is about feedback rather than hierarchy. Executing a precomputed plan to fixed end points fails as soon as
the supports move appreciably: no policy completes any of 30 trials at 10\,cm amplitude. With continuous
tracking of the moving target, all five feedback policies complete 30 of 30 trials in every moving condition,
while the open-loop policy completes none of 120. The cost of motion under ideal tracking is small (about
1\,\% more construction time for local feedback at 10\,cm and 0.05\,Hz). The same study records what did
\emph{not} help: predictive repair and frequent full replanning add planning work without improving
completion (claims C07, C08), the HTN agent and a flat agent given the same sequence behave identically
(C06), and completion collapses when sensing lags by one second or more (C10).

\paragraph{Recovery from execution failure.}
The recovery benchmark (Section~\ref{sec:recovery}) separates search backtracking from recovery during
execution and tests the latter with the re-entrant planner IPyHOP. All four back-ends complete every trial
in the five feasible failure scenarios and correctly report the two infeasible ones. Retaining the solution
tree reduces method calls in most scenarios but increases them when silk runs out, and it does not reduce
planning time as implemented (claims C21, C26). Integrated into the moving-support model, retained and fresh
IPyHOP give identical physical outcomes in all 720 paired trials (C25).

\paragraph{A physical picture of the planned web.}
The simulation of Section~\ref{sec:sim} adds what the planner's straight-line drawings leave out: threads
that hang under gravity, belly in wind and deform with moving supports. Its cable solver reproduces the exact elastic catenary to within 0.02\,\% of the sag on a
16\,cm span. With the illustrative silk values used, such a span sags only 0.037\,cm, so the figures
multiply the silk's weight by a stated factor to make the curvature visible. The simulation is an
illustration layered on the planner's output and feeds no claim in the ledger.

\paragraph{The research record.}
Since October the project keeps a dated work log, a claim ledger and hashed evidence files
(Sections~\ref{sec:record} and~\ref{sec:record-b}). Of the 28 claims the ledger tracks, 11 are supported, 10
were tested and found negative, 6 are not established and 1 is explicitly not claimed. All 177 tests pass.

\subsection{Discussion}
\label{sec:synthesis-discussion}

Three points follow from taking these results together.

First, what the HTN encodes is knowledge of the construction order, and that is what the comparison with a
random walker measures. The tenfold energy gap is real and stable, but it separates an agent that knows the
order from one that does not. Where a flat agent was handed the same action sequence, in the moving-support
benchmark, it performed identically to the HTN agent. No experiment in the repository yet shows a benefit
that comes from the hierarchy as such.

Second, robustness in these models comes from closing the loop during execution. Fixed plans fail under
moving supports whatever planner produced them, and simple local feedback is sufficient under ideal sensing.
More elaborate replanning has so far bought nothing measurable, and the benefit of retaining a solution tree
is confined to planning effort in particular failure scenarios.

Third, the biological reading must stay modest. The literature review (Section~\ref{sec:record-related})
supports describing web building in stages with cue-dependent behaviour, which is what an HTN can express,
but it does not establish that spiders plan hierarchically, and the simulator is not calibrated against
measured construction data (claim C12). The link to a caloric optimisation proposed at the end of
Section~\ref{sec:orig-discussion} remains a proposal: the pruning heuristic it relies on is not implemented.

\subsection{Issues found while preparing this update}
\label{sec:synthesis-issues}

Checking every number against the repository surfaced the following defects. None is fixed in commit
\texttt{61d18d8}; each is small and would make the affected result reproducible or comparable.

\begin{enumerate}
\item The A* heuristic passes coordinates to a distance function that expects node names, so it is constant
or zero rather than informative, and it is not admissible even when corrected because two rig nodes share a
position (Section~\ref{sec:astar_fap-effort}).
\item The random-walker baseline enumerates candidate actions from a Python set, so its results depend on
the interpreter's hash seed and the committed baseline rows cannot be regenerated exactly
(Section~\ref{sec:core-anchor}).
\item The HTN agent and the baseline are scored by different completion criteria in the sensitivity sweeps
(Section~\ref{sec:core-anchor}).
\item No runner of the original or updated HTN model selects a GTPyhop search strategy, so the documented
backtracking is not what executes (Section~\ref{sec:recovery-strategy}).
\item The spiral generator steps through seven directions while the configuration assumes eight per turn, so
the spirals make more rotations than configured and one radius receives no spiral attachment
(Section~\ref{sec:core-spiral}).
\item The \texttt{empirical:} configuration block and \texttt{empirical\_data.md} are not read by any code
(Sections~\ref{sec:core-config}, \ref{sec:astar_fap-empirical}).
\end{enumerate}

\subsection{Conclusions}
\label{sec:synthesis-conclusions}

The project now has an executable and configurable HTN model of orb-web construction after Zschokke, two
alternative formulations, a moving-support benchmark with six execution policies, a recovery benchmark built
on a re-entrant planner, and a physical simulation of the planned web. The formalisation of Checkpoint~2
stands. Its comparative claims are narrowed: the energy advantage over undirected search is confirmed, the
reliability advantage is withdrawn, and backtracking over alternative strategies is encoded but has not been
exercised. The new experiments show that closed-loop execution, not hierarchical structure, is what carries
construction through disturbance in these models, and they record the negative results as carefully as the
positive ones. The open questions and the experiments proposed to address them are listed in
Section~\ref{sec:record-limits}.
